\documentclass[10pt,a4paper]{amsart}
\usepackage[margin=19mm]{geometry}
\usepackage{amsmath,amssymb,mathtools}
\usepackage[scr=rsfs,cal=euler]{mathalfa}
\usepackage{xfrac}
\usepackage[unicode,hidelinks,bookmarks=false]{hyperref}
\newcommand{\ie}{i.e.\ }
\newcommand{\cf}{cf.\ }
\newcommand{\resp}{respectively\ }
\newcommand{\Subsection}{\subsection}
\providecommand{\nobreakdash}{\nobreak\hbox{-}}
\setlength{\emergencystretch}{2em}
\multlinegap0pt
\usepackage{bm}
\usepackage{bbm}
\usepackage{dsfont}
\usepackage{xspace}
\usepackage{cancel}
\usepackage{mathtools}
\usepackage{amsthm}
\makeatletter
\@ifundefined{theorem}{\newtheorem{theorem}{Theorem}}{}
\@ifundefined{proposition}{\newtheorem{proposition}[theorem]{Proposition}}{}
\makeatother
\newcommand{\R}{\mathbb R}
\title[Log-concavity of solutions of parabolic equations related to]{Log-concavity of solutions of parabolic equations related to the Ornstein-Uhlenbeck operator and applications}
\author{Andrea Colesanti, Lei Qin, Paolo Salani}
\date{Source: arXiv:2601.07426v2 (CC BY 4.0)}
\begin{document}
\maketitle

\noindent\emph{Source and licence.} Log-concavity of solutions of parabolic equations related to the Ornstein-Uhlenbeck operator and applications. Version arXiv:2601.07426v2 (CC BY 4.0), distributed under the Creative Commons Attribution 4.0 International licence, as recorded in the arXiv licence metadata for this exact version and shown on the abstract page https://arxiv.org/abs/2601.07426v2. Full rights notice: licenses/arxiv-2601.07426-CC-BY-4.0.txt.\par\smallskip

\noindent\textit{Editorial scope.} Counted unit: the Proposition whose source label is \texttt{None} in the article (source0.tex lines 469--475, source label \texttt{None}), delivered together with the article's own complete original proof of it (source0.tex lines 477--493, one \texttt{proof} environment of 17 lines). No mathematics is written, repaired or re-derived: statement and proof are the article's own text line for line. Required context retained uncounted: parabolic-equation1 (lines 134--142); heat-kernel-1 (lines 323--325); log-concavity proposition (lines 337--346); H-PDE1 (lines 458--465). Every definition, display and label of those lines is retained unchanged. Omitted from this package and not redistributed: 1--133, 143--322, 326--336, 347--457, 466--468, 476--476, 494--971. Nothing inside a retained excerpt is omitted or altered: every retained body is delivered exactly as the article's lines are written, so no omission receipt of any kind accompanies this package. Citations: the retained excerpts carry no citation command at all, so no bibliography entry of the article is transcribed and no reference is redistributed. Reference adaptation: none was needed and none was made; every reference inside the delivered unit resolves inside it, so no unresolved reference or citation is produced. Printed slips kept literal: no printed word, symbol, hypothesis or number of the counted statement or of its proof was altered. Layout and class: the article's own class is \texttt{amsart} with packages including geometry, amsmath, amssymb, amscd, amsthm, latexsym, graphicx, babel, inputenc, textcomp, times, colortbl; the wrapper is the lane's pinned \texttt{amsart} editorial wrapper at 10pt on A4, which restates the theorem-like environments the retained text needs and adds the article's own macro definitions that the retained text needs (\texttt{\textbackslash R}), with extra wrapper packages (bm, bbm, dsfont, xspace, cancel, mathtools). The delivered numbering is the unit's own. Assets: no retained span contains a figure environment, a graphics-inclusion command, a PDF-page inclusion, a table or a TikZ picture; no image file is copied into this package, the package declares no asset dependency and no image file is redistributed with it. Licence: version arXiv:2601.07426v2 of this article is distributed under the Creative Commons Attribution 4.0 International licence, as recorded in the arXiv licence metadata for this exact version and shown on the abstract page https://arxiv.org/abs/2601.07426v2. A full rights notice is delivered at licenses/arxiv-2601.07426-CC-BY-4.0.txt. Characters: every retained excerpt is pure ASCII, so the delivered engine, XeLaTeX with its default Latin Modern text font, reports no missing character.
\begin{equation}\label{parabolic-equation1}
\left\{
\begin{aligned}
& \Delta_x p_\gamma^\Omega(x,z;t)-(x,\nabla_x p_\gamma^\Omega) =\partial_t p_\gamma^\Omega(x,z;t), \quad\forall\ x,z\in\Omega, \ t>0,  \\
& p_\gamma^\Omega(x,z;t)=0, \quad\mbox{if $x\in\partial\Omega$ or $z\in\partial\Omega$, $t>0$},  \\
& p_\gamma^\Omega(x,z;0)d\gamma(z)=\delta(x-z),  \quad  \forall\ x,z\in\Omega,
\end{aligned}
\right.
\end{equation}

\begin{equation}\label{heat-kernel-1}
p_\gamma(x,z;t):= \frac{1}{(1-e^{-2t})^{n/2}} \text{exp} \left( -\frac{|x|^2-2e^t (x,z)+|z|^2}{2(e^{2t}-1)}\right).
\end{equation}

\begin{proposition}\label{log-concavity proposition} For every $t>0$, the function
$$
(x,z)\, \rightarrow\,p_\gamma (x,z;t) e^{-|z|^2/2},\quad(x,z)\in\R^n\times\R^n
$$
is log-concave. Moreover, if $g\colon\R^n\to\R$ is log-concave, then, for every $t>0$, the function $G\colon\R^n\to\R$ defined by
$$
G(x)=\int_{\R^n} g(z)p_\gamma(x,z;t)d\gamma(z),
$$
is log-concave.
\end{proposition}

\begin{equation}\label{H-PDE1}
\left\{
\begin{aligned}
& \Delta u -(x,\nabla u)=\partial_t u, \quad & \mbox{in $\R^n\times \{t>0\}$,}\\
& u(x,0)=u_0(x),  \quad & \mbox{on $\R^n \times \{ t=0 \}$.}
\end{aligned}
\right.
\end{equation}

\begin{proposition}
Let $u_0\in L^2(\R^n,\gamma)$. Then the function $u\colon\R^n\times[0,\infty)\to\R$, defined by
\[
u(x,t)=\int_{\R^n} u_0(e^{-t}x+\sqrt{1-e^{-2t}}y) d\gamma(y),
\]
is a solution of problem \eqref{H-PDE1}. Furthermore, if $u_0$ is log-concave in $\R^n$, then $u(\cdot,t)$ is log-concave for every $t>0$.
\end{proposition}

\begin{proof}[\bf Proof]
By \eqref{heat-kernel-1}, we can directly check that the kernel $p_\gamma$ solves the parabolic equation in problem \eqref{parabolic-equation1}. As a consequence, the function $u$ defined by
\begin{equation}\label{Pt}
\begin{aligned}
u(x,t)=\int_{\R^n} u_0(z) p_\gamma(x,z;t) d\gamma(z)
\end{aligned}
\end{equation}
is a solution of the PDE in problem \eqref{H-PDE1}. Moreover, by the change of variable $z=e^{-t}x+\sqrt{1-e^{-2t}}y$, we obtain the following expression for $u$:
\begin{equation}\nonumber
\begin{aligned}
u(x,t)=\int_{\R^n} u_0(e^{-t}x+\sqrt{1-e^{-2t}}y) d\gamma(y).
\end{aligned}
\end{equation}
Clearly, letting $t=0$, we have $u(x,t)=u_0$. The previous facts show that $u$ is a solution of the equation in \eqref{H-PDE1}.

Finally, if $u_0$ is log-concave, then the log-concavity of $u(\cdot,t)$, for every $t>0$, follows from \eqref{Pt} and Proposition \ref{log-concavity proposition}.
\end{proof}

\end{document}
